\documentclass[11pt]{amsart}

\usepackage[margin=1in]{geometry}
\usepackage{amsmath,amssymb,amsthm,mathtools}
\usepackage[shortlabels]{enumitem}
\usepackage{hyperref}
\usepackage{setspace}

\numberwithin{equation}{section}

\newtheorem{theorem}{Theorem}[section]
\newtheorem{proposition}[theorem]{Proposition}
\newtheorem{lemma}[theorem]{Lemma}
\newtheorem{corollary}[theorem]{Corollary}
\theoremstyle{definition}
\newtheorem{definition}[theorem]{Definition}
\theoremstyle{remark}
\newtheorem{remark}[theorem]{Remark}
\newtheorem{question}[theorem]{Question}

\newcommand{\E}{\mathbb E}
\newcommand{\Prob}{\mathbb P}
\newcommand{\N}{\mathbb N}
\newcommand{\R}{\mathbb R}
\newcommand{\Span}{\operatorname{span}}
\newcommand{\supp}{\operatorname{supp}}
\newcommand{\sgn}{\operatorname{sign}}
\newcommand{\conc}{\operatorname{Q}}
\newcommand{\eps}{\varepsilon}
\newcommand{\abs}[1]{\lvert #1 \rvert}
\newcommand{\norm}[1]{\lVert #1 \rVert}
\newcommand{\ip}[2]{\langle #1,#2\rangle}

\title[Stable Phase Retrieval for Independent Coordinates]{Stable Phase Retrieval for Spans of Independent Random Variables}

\author{Pedro Abdalla Teixeira\textsuperscript{1}}
\author{Jaume de Dios Pont\textsuperscript{2}}
\author{Jo\~ao Pedro Ramos\textsuperscript{3}}
\author{Mitchell Taylor\textsuperscript{4}}

\begin{document}
\setstretch{1.2}

\begin{abstract}
We study stable phase retrieval for closed spans of independent real-valued
random variables in $L^2$.  We prove the sufficiency direction of the
Calderbank--Daubechies--Freeman--Freeman conjecture up to one exceptional
coordinate: after $L^2$ normalization and relabeling, stable phase retrieval
holds whenever all but possibly one coordinate satisfy a uniform two-sided
$L^1$ bound.  The paper gives two proofs.  The first is a compactness proof
based on a head--tail decomposition, bounded probe functions, and infinitely
divisible limits.  The second is a quantitative proof in the homogeneous case,
using an explicit dichotomy and anticoncentration estimates for diffuse tails.
\end{abstract}

\maketitle
\begingroup
\renewcommand{\thefootnote}{\arabic{footnote}}
\footnotetext[1]{Department of Mathematics, University of California, Irvine, Irvine, California 92697, USA}
\footnotetext[2]{Center for Data Science, New York University, New York, New York 10011, USA}
\footnotetext[3]{Instituto de Matem\'atica Pura e Aplicada (IMPA), Rio de Janeiro, RJ, Brazil}
\footnotetext[4]{Department of Mathematics, ETH Z\"urich, Z\"urich, Switzerland}
\endgroup
\setcounter{footnote}{0}

\section{Introduction}

The phase retrieval problem asks whether an unknown vector or function can be
recovered from phaseless data.  In the real setting, the unavoidable symmetry
is multiplication by a global sign, and the central issue is whether this is
the only obstruction.  In finite dimensions, phase retrieval has by now a
substantial theory, spanning injectivity, stability, convex relaxations,
random measurements, and structured models
\cite{bandeira2014saving,eldar2014stability,MR3069958,MR3260258,KS,DB}.
Infinite-dimensional phase retrieval is subtler.  Recovery itself can hold in
Hilbert spaces \cite{cahill2016infinite}, but uniform stability can fail for
continuous frames \cite{alaifari2017continuous}, and natural measurement
systems such as Gabor families can be severely ill-posed
\cite{alaifari2021gabor,grohs2021stable,grohs20212,steinerberger2020stability}.
For general background we refer to \cite{grohs2020survey}.
Further landmarks include early reconstruction and finite-dimensional stability
results
\cite{sanz1984phase,ramos-sousa-pauli,balan2013,bandeira2014saving,eldar2014stability,MR3069958,MR3260258,KS,DB,MR3746047},
variants for affine, vector-valued, and projection-based models
\cite{bartusel2021phase,chen2019phase,casazza2013,casazza2017weak,casazza2017norm,akrami2021note},
and infinite-dimensional or structured-measurement results
\cite{calderbank2022stable,localphase,MR4298599}.
For the surrounding frame-theoretic and discretization background we also refer
to
\cite{AG-continuous-frames,CHL,CCS-JMAA,CDOS,FG-atomic,MR3275625,MR3367650,MR3901674,fg22,kashin2021sampling,LT21,DPSTT,NOU,MR2011364,MR0423039,luef2021gaussian}.

Two themes run through the subject.  The first is uniqueness: does the map
$f\mapsto |f|$ distinguish points of the model space modulo the natural
symmetry?  The second is stability: can one recover the signal in a uniformly
continuous, or even Lipschitz, fashion from phaseless data?  In applications,
the second question is indispensable.  Exact injectivity without stability is
too weak for numerical reconstruction, and many infinite-dimensional models
exhibit precisely this pathology.  One of the main conceptual contributions of
\cite{FOTP} is to make the stability question intrinsic to the geometry of the
subspace rather than to a particular discretized frame.  This point of view is
particularly well suited to random subspaces of $L^2$, where the ambient
structure is probabilistic from the outset and where orthogonality interacts
naturally with independence.  In a complementary direction, Christ, Pineau,
and Taylor \cite{christ2022examples} constructed examples with only
H\"older-stable recovery, underscoring that even when phase retrieval remains
possible the quantitative modulus can deteriorate substantially.

The model considered here is also connected to the geometry of independent
coordinates in classical Banach space theory.  An independent sequence of
centered $L^2$-normalized random variables behaves like a random coordinate
system whose linear span carries both Euclidean and probabilistic structure.
The size of the $L^1$ norm of each coordinate is a first quantitative measure
of how much sign information survives after passing to absolute values.
Variables with very small $L^1$ norm are too concentrated near zero, while
variables with $L^1$ norm too close to one are nearly flat in modulus and
hence almost insensitive to sign changes.  The theorem proved in this paper
says, in essence, that outside of those two obstructions there is enough
randomness in the coordinate system to force uniform phase retrieval.

It is worth stressing that the problem is not coordinatewise.  Stable phase
retrieval asks for control of \emph{all} vectors in the closed span, not merely
for the generators.  Thus one must understand how arbitrary linear
combinations behave under the nonlinear map $f\mapsto |f|$.  The difficult
regime is when two orthogonal combinations have almost identical moduli.  In
that situation the coefficients may spread over infinitely many coordinates,
and no finite-dimensional argument can be applied directly.  The
head--tail decomposition used below is precisely the mechanism that overcomes
this: the few large coordinates form the head, which is rigid enough to be
identified by bounded probes, while the many small coordinates form the tail,
which is diffuse enough to admit probabilistic structure theorems.

We adopt the function-space viewpoint introduced in
\cite{FOTP,AG-continuous-frames}.  Let $X$ be a real function space and let
$E\subset X$ be a subspace.  We say that $E$ does \emph{stable phase
retrieval} if there exists $C<\infty$ such that
\[
\min_{\sigma\in\{\pm1\}}\norm{f-\sigma g}_X
\le C\,\norm{\abs{f}-\abs{g}}_X
\qquad\text{for all }f,g\in E.
\]
The theorem of Freeman, Oikhberg, Pineau, and Taylor \cite{FOTP} shows that in
real Hilbert spaces the problem reduces to orthogonal pairs: to prove stable
phase retrieval on a subspace it suffices to lower-bound the modulus gap for
pairs of orthogonal unit vectors in that subspace.  This orthogonal reduction
is the common point of departure of both proofs given here.  For historical
continuity, we note that \cite{our-paper} is the preprint version of
\cite{FOTP}.

Our interest is the structure of subspaces generated by independent random
variables.  Calderbank, Daubechies, Freeman, and Freeman
\cite{calderbank2022stable} proved stable phase retrieval for a shifted model
and proposed the following conjectural characterization in the unshifted
setting: if $(\xi_i)$ is an orthonormal sequence of independent centered
random variables in $L^2$, then $\overline{\Span}(\xi_i)$ does stable phase
retrieval if and only if there exist constants $0<A\le B<1$ such that
$A\le \norm{\xi_i}_{L^1}\le B$ for all but at most one index.  The lower bound
excludes concentration near zero; the upper bound excludes the Rademacher
regime in which $\abs{\xi_i}$ is constant and therefore carries no sign
information.

The phrase ``at most one exceptional coordinate'' is not an accident of the
proof.  It is the largest defect compatible with the expected geometry of the
problem.  One exceptional coordinate may coexist with an otherwise
phase-retrieving family, but two genuinely exceptional coordinates already
allow a nontrivial loss of sign information at the level of the span.  The
one-exception theorem should therefore be viewed as the natural endpoint of the
qualitative theory.

The necessity of these conditions is by now standard.  The equivalence of
$L^1$ and $L^2$ norms on a phase-retrieving subspace follows from
\cite{FOTP}; in particular a uniform positive lower $L^1$ bound is forced.
Likewise, a centered $L^2$-normalized variable satisfies
$\norm{\xi}_{L^1}=1$ if and only if $\xi^2\equiv 1$, that is, if and only if
$\xi$ is Rademacher-valued.  But a span containing two Rademacher directions
cannot even do phase retrieval, since all such coordinates have the same
absolute value.

Our first theorem gives the desired sufficiency result with one exceptional
coordinate.

\begin{theorem}[One-exception stable phase retrieval]\label{thm:main}
Let $\eta,\xi_2,\xi_3,\dots$ be independent real-valued random variables and
fix $\delta>0$.  Assume
\[
\norm{\eta}_{L^2}=1,
\qquad
\E[\eta]=0,
\]
and, for every $i\ge2$,
\[
\norm{\xi_i}_{L^2}=1,
\qquad
\E[\xi_i]=0,
\qquad
\delta\le \norm{\xi_i}_{L^1}\le 1-\delta.
\]
Then there exists $C_{\delta,\eta}\ge1$ such that, for all
$X,Y\in \overline{\Span}(\eta,\xi_i:i\ge2)$,
\[
\min_{\sigma\in\{\pm1\}}\norm{X-\sigma Y}_{L^2}
\le
C_{\delta,\eta}\,\norm{\abs{X}-\abs{Y}}_{L^2}.
\]
\end{theorem}

\begin{remark}
After relabeling the coordinates and normalizing in $L^2$,
Theorem~\ref{thm:main} is exactly the sufficiency direction of the
Calderbank--Daubechies--Freeman--Freeman conjecture up to one exceptional
variable.
\end{remark}

The second proof is quantitative in spirit.  It keeps the one-exception
framework of Theorem~\ref{thm:main}, but replaces the limit-law endpoint by an
explicit head--tail dichotomy and by anticoncentration estimates for the
diffuse part.  In that sense, the compactness and quantitative proofs address
the same theorem from complementary angles.  The compactness proof shows that a
bad sequence cannot exist by passing to a limit model; the quantitative proof
shows directly that every orthogonal pair already contains a definite amount of
modulus separation.

The two proofs illuminate complementary aspects of the theorem.  The
compactness proof begins with a normalized counterexample array and extracts a
head--tail decomposition of the coefficients.  The head consists of the few
coordinates carrying macroscopic mass, while the tail is diffuse in the sense
that its maximal coefficient goes to zero.  Bounded probes read off quadratic
coefficients from the observable $|X_n|^2-|Y_n|^2$, forcing the head limit to
agree up to sign.  The tail then converges to an infinitely divisible law.
Since the exact limit identity still has the form
$|H_X+R_X|=|H_Y+R_Y|$, the support of the tail law is forced onto one of the
phase-retrieval lines; tail rigidity then contradicts the orthogonal
normalization.  In the degenerate-head regime, the same contradiction can be
reached by a different endpoint argument: diffuse tails produce a nontrivial
Gaussian component in the limit law, and a line-geometry argument for
convolutions supported on $\{(u,v):\abs{u}=\abs{v}\}$ forces that Gaussian
component, and hence the full law, onto a single line.

The quantitative proof replaces the limit law by an explicit dichotomy.  The
coordinates are split into a head and a tail according to a fixed cutoff.
Either one of the signs $a\pm b$ is essentially supported on the head, in
which case bounded probes recover the head difference matrix and orthogonality
forces a positive modulus gap, or both $a+b$ and $a-b$ remain diffuse on the
tail.  In the latter regime one partitions the tail into many blocks of
comparable energy and applies symmetrization, Sperner-type combinatorics, and
small-ball estimates to prove that both combinations stay away from zero with
positive probability.  This yields an explicit lower bound for
$\E\abs{X_a^2-X_b^2}$ and hence for $\norm{\abs{X_a}-\abs{X_b}}_{L^2}$.

The paper is organized as follows.  Section~\ref{sec:preliminaries} collects
the basic definitions and the analytic tools used throughout.  In
Section~\ref{sec:compactness-proof} we present the compactness proof of
Theorem~\ref{thm:main}, including the Gaussian alternative to the final
line-support argument.  Section~\ref{sec:quant-proof} gives a second proof of
Theorem~\ref{thm:main}, explaining how the limit-law step can be bypassed by an
explicit decomposition together with anticoncentration.
Finally, Section~\ref{sec:comments} discusses extensions, related directions,
and open problems.  Although we now have a quantitative proof in the
homogeneous setting, we have chosen to retain the compactness proof because of
its flexibility: it isolates a robust mechanism that is likely to survive in
more general symmetry classes and in settings where quantitative control is
secondary.

\section{Preliminaries}\label{sec:preliminaries}

\subsection{Stable phase retrieval and orthogonal reduction}

\begin{definition}
Let $E$ be a subspace of a real Banach lattice $X$.  We say that $E$ does
\emph{stable phase retrieval} if there exists $C<\infty$ such that
\[
\min_{\sigma\in\{\pm1\}}\norm{f-\sigma g}_X
\le
C\,\norm{\abs{f}-\abs{g}}_X
\qquad(f,g\in E).
\]
\end{definition}

We shall use repeatedly the following orthogonal reduction from \cite{FOTP}.

\begin{proposition}[Orthogonal reduction]\label{prop:orth-reduction}
Let $H$ be a real Hilbert space and let $E\subset H$ be a subspace.  Then $E$
does stable phase retrieval if and only if there exists $c>0$ such that
\[
\norm{\abs{f}-\abs{g}}_H\ge c
\]
for every orthogonal pair $f,g\in E$ with $\norm{f}_H=\norm{g}_H=1$.
\end{proposition}

\subsection{Probe functions and lower $L^1$ bounds}

The first input is the existence of bounded dual probes.  These are the basic
coefficient-extraction devices in both proofs.

\begin{lemma}[Linear and quadratic probes]\label{lem:probes}
Let $\xi$ be a real random variable with $\E[\xi]=0$ and $\E[\xi^2]=1$.

\begin{enumerate}[label=\rm(\arabic*)]
\item If $\E[\abs{\xi}]\ge A>0$, then there exists a bounded measurable
function $S$ such that
\[
\E[S]=0,
\qquad
\E[\xi S]=1,
\qquad
\norm{S}_{L^\infty}\le \frac{2}{A}.
\]
\item If $\E[\abs{\xi}]\le B<1$, then there exists a bounded measurable
function $T$ such that
\[
\E[T]=0,
\qquad
\E[(\xi^2-1)T]=1,
\qquad
\norm{T}_{L^\infty}\le \frac{2}{1-B}.
\]
\end{enumerate}
\end{lemma}

\begin{proof}
We first construct the linear probe.  Define
\[
S_0:=\sgn(\xi),
\]
with the convention $\sgn(0)=0$.  Then $S_0$ is measurable and satisfies
$|S_0|\le 1$ almost surely.  Moreover,
\[
\xi S_0=\abs{\xi}
\qquad\text{almost surely,}
\]
so
\[
\E[\xi S_0]=\E[\abs{\xi}]\ge A>0.
\]
This positivity allows us to normalize.  Set
\[
S:=\frac{S_0-\E[S_0]}{\E[\xi S_0]}.
\]
The function $S$ is measurable and bounded.  Its mean is
\[
\E[S]=\frac{\E[S_0]-\E[S_0]}{\E[\xi S_0]}=0,
\]
and, using that $\E[\xi]=0$,
\[
\E[\xi S]
=\frac{\E[\xi S_0]-\E[S_0]\E[\xi]}{\E[\xi S_0]}
=\frac{\E[\xi S_0]}{\E[\xi S_0]}
=1.
\]
Finally, since $|S_0|\le1$ and $|\E[S_0]|\le \E|S_0|\le1$,
\[
|S_0-\E[S_0]|\le 2
\qquad\text{almost surely,}
\]
and therefore
\[
\norm{S}_{L^\infty}
\le \frac{2}{\E[\xi S_0]}
=\frac{2}{\E[\abs{\xi}]}
\le \frac{2}{A}.
\]
This proves part~\rm(1).

For the quadratic probe, apply the same idea to the centered variable
\[
\eta:=\xi^2-1.
\]
Because $\E[\xi^2]=1$, one has $\E[\eta]=0$.  We now show that $\eta$ has
positive $L^1$ mass bounded away from zero.  Since $\xi^2=\abs{\xi}^2$,
\[
\abs{\eta}
=\abs{\abs{\xi}^2-1}
=\abs{\abs{\xi}-1}(\abs{\xi}+1)
\ge \abs{\abs{\xi}-1}.
\]
Taking expectations and using the reverse triangle inequality yields
\[
\E[\abs{\eta}]
\ge \E[\abs{\abs{\xi}-1}]
\ge \bigl|\E[\abs{\xi}]-1\bigr|
=1-\E[\abs{\xi}]
\ge 1-B.
\]
Now set
\[
T_0:=\sgn(\eta)
\]
and define
\[
T:=\frac{T_0-\E[T_0]}{\E[\eta T_0]}
=\frac{T_0-\E[T_0]}{\E[\abs{\eta}]}.
\]
Exactly as in the linear case, $T$ is measurable, bounded, and satisfies
\[
\E[T]=0,
\qquad
\E[\eta T]=1.
\]
Since $\eta=\xi^2-1$, this reads
\[
\E[(\xi^2-1)T]=1.
\]
The $L^\infty$ bound is
\[
\norm{T}_{L^\infty}
\le \frac{2}{\E[\abs{\eta}]}
\le \frac{2}{1-B}.
\]
This completes the proof.
\end{proof}

\begin{remark}\label{rem:probe-extraction}
If $\xi_1,\dots,\xi_n$ are independent centered random variables and
$X_c=\sum_i c_i\xi_i$, then the centering of the probes eliminates every
coordinate except the one being tested.  In particular,
\[
\E[X_c\,S_i]=c_i,
\qquad
\E[X_c^2\,T_i]=c_i^2,
\qquad
\E[X_c^2\,S_iS_j]=2c_ic_j
\quad(i\ne j),
\]
whenever $S_i,S_j,T_i$ are the corresponding probes for $\xi_i,\xi_j$.
\end{remark}

The second basic estimate is the lower $L^1$ bound for independent sums.

\begin{lemma}[$L^1$ lower bound for independent sums]\label{lem:l1lower}
Let $\zeta_1,\dots,\zeta_N$ be independent centered real random variables and
assume that
\[
\norm{\zeta_i}_{L^1}\ge \delta \norm{\zeta_i}_{L^2}
\qquad(1\le i\le N).
\]
Then for every scalar family $(c_i)_{i=1}^N$,
\[
\norm{\sum_{i=1}^N c_i\zeta_i}_{L^1}
\ge
\frac{\delta}{2\sqrt2}
\Bigl(\sum_{i=1}^N c_i^2\norm{\zeta_i}_{L^2}^2\Bigr)^{1/2}.
\]
\end{lemma}

\begin{proof}
The statement is homogeneous in the coefficient vector, so we may first
normalize and assume that
\[
\sum_{i=1}^N c_i^2\norm{\zeta_i}_{L^2}^2=1.
\]
Indeed, once the claim is established under this normalization, the general
case follows by scaling.

We next reduce to the unit-variance case.  For each $i$ set
\[
\widetilde\zeta_i:=\frac{\zeta_i}{\norm{\zeta_i}_{L^2}},
\qquad
\widetilde c_i:=c_i\norm{\zeta_i}_{L^2}.
\]
Then the variables $\widetilde\zeta_i$ are still independent and centered,
they satisfy $\norm{\widetilde\zeta_i}_{L^2}=1$, and the lower $L^1$ bound
becomes
\[
\norm{\widetilde\zeta_i}_{L^1}
=\frac{\norm{\zeta_i}_{L^1}}{\norm{\zeta_i}_{L^2}}
\ge \delta.
\]
Moreover,
\[
\sum_{i=1}^N \widetilde c_i \widetilde\zeta_i
=\sum_{i=1}^N c_i\zeta_i,
\qquad
\sum_{i=1}^N \widetilde c_i^2
=1.
\]
Thus it is enough to prove the lemma under the additional assumptions
\[
\norm{\zeta_i}_{L^2}=1
\qquad(1\le i\le N),
\qquad
\sum_{i=1}^N c_i^2=1.
\]

Define
\[
S:=\sum_{i=1}^N c_i\zeta_i,
\]
and let
\[
S':=\sum_{i=1}^N c_i\zeta_i'
\]
be an independent copy of $S$, obtained from an independent copy
$(\zeta_i')_{i=1}^N$ of the family $(\zeta_i)_{i=1}^N$.
By the triangle inequality,
\[
\abs{S-S'}\le \abs{S}+\abs{S'}.
\]
Taking expectations and using that $S$ and $S'$ have the same distribution, we
obtain
\[
\norm{S-S'}_{L^1}\le 2\norm{S}_{L^1}.
\]
Equivalently,
\[
2\norm{S}_{L^1}\ge \norm{S-S'}_{L^1}.
\]

Now the random variables $\zeta_i-\zeta_i'$ are independent and symmetric.
Conditioning on their absolute values and applying the Khintchine inequality
gives
\[
\norm{S-S'}_{L^1}
=
\Bigl\|\sum_{i=1}^N c_i(\zeta_i-\zeta_i')\Bigr\|_{L^1}
\ge
\frac{1}{\sqrt2}\,
\E\Bigl[\Bigl(\sum_{i=1}^N c_i^2(\zeta_i-\zeta_i')^2\Bigr)^{1/2}\Bigr].
\]

Because the nonnegative numbers $(c_i^2)_{i=1}^N$ sum to one, Cauchy--Schwarz
implies the pointwise estimate
\[
\sum_{i=1}^N c_i^2\abs{\zeta_i-\zeta_i'}
\le
\Bigl(\sum_{i=1}^N c_i^2\Bigr)^{1/2}
\Bigl(\sum_{i=1}^N c_i^2(\zeta_i-\zeta_i')^2\Bigr)^{1/2}
=
\Bigl(\sum_{i=1}^N c_i^2(\zeta_i-\zeta_i')^2\Bigr)^{1/2}.
\]
Taking expectations and combining with the previous inequality, we arrive at
\[
\norm{S-S'}_{L^1}
\ge
\frac{1}{\sqrt2}\sum_{i=1}^N c_i^2\norm{\zeta_i-\zeta_i'}_{L^1}.
\]
Hence
\[
\norm{S}_{L^1}
\ge
\frac{1}{2\sqrt2}\sum_{i=1}^N c_i^2\norm{\zeta_i-\zeta_i'}_{L^1}.
\]

It remains to bound each $\norm{\zeta_i-\zeta_i'}_{L^1}$ from below.  Since
$\zeta_i'$ is independent of $\zeta_i$ and centered,
\[
\E[\zeta_i-\zeta_i' \mid \zeta_i]
=\zeta_i.
\]
Applying Jensen's inequality to the conditional expectation gives
\[
\norm{\zeta_i-\zeta_i'}_{L^1}
\ge
\norm{\E[\zeta_i-\zeta_i' \mid \zeta_i]}_{L^1}
=\norm{\zeta_i}_{L^1}
\ge \delta.
\]
Substituting this into the previous bound yields
\[
\norm{S}_{L^1}
\ge
\frac{\delta}{2\sqrt2}\sum_{i=1}^N c_i^2
=\frac{\delta}{2\sqrt2}.
\]
Since this is exactly the normalized form of the desired estimate, the proof is
complete.
\end{proof}

\subsection{Diffusive arrays and infinitely divisible limits}

We shall use the classical fact that weak limits of infinitesimal triangular
arrays of independent random variables are infinitely divisible; see, for
instance, \cite[Chapter~IV]{gnedenko-kolmogorov}.

\begin{lemma}[Diffuse tail rigidity]\label{lem:tail-rigidity}
Let $(\xi_{n,i})_{n\ge1,\ i\ge2}$ be a triangular array such that for each
fixed $n$ the family $(\xi_{n,i})_{i\ge2}$ is independent, centered,
$L^2$-normalized, and satisfies
\[
\norm{\xi_{n,i}}_{L^1}\ge \delta
\qquad(i\ge2).
\]
Let $c^{(n)}\in \ell^2(\N)$ satisfy
\[
\sup_n \norm{c^{(n)}}_2<\infty,
\qquad
c_1^{(n)}=0,
\qquad
\max_{i\ge2}\abs{c_i^{(n)}}\to0.
\]
Set
\[
Z_n:=\sum_{i\ge2} c_i^{(n)}\xi_{n,i}.
\]
If $Z_n\to0$ in distribution, then $\norm{c^{(n)}}_2\to0$.
\end{lemma}

\begin{proof}
Because the limiting law is the Dirac mass at $0$, the assumption
$Z_n\to0$ in distribution implies that
\[
Z_n\to0
\qquad\text{in probability.}
\]
To pass from convergence in probability to convergence of first moments, we use
uniform integrability.

For each $n$, the variables $(\xi_{n,i})_{i\ge2}$ are independent, centered,
and $L^2$-normalized.  Therefore they form an orthonormal family in $L^2$, and
the random sum $Z_n$ has
\[
\norm{Z_n}_{L^2}^2
=\E\Bigl[\Bigl(\sum_{i\ge2} c_i^{(n)}\xi_{n,i}\Bigr)^2\Bigr]
=\sum_{i\ge2}(c_i^{(n)})^2\E[\xi_{n,i}^2]
=\sum_{i\ge2}(c_i^{(n)})^2
=\norm{c^{(n)}}_2^2.
\]
Since $\sup_n \norm{c^{(n)}}_2<\infty$, the sequence $(Z_n)_n$ is bounded in
$L^2$.  Any $L^2$-bounded family is uniformly integrable in $L^1$, so
$(|Z_n|)_n$ is uniformly integrable.  Because $Z_n\to0$ in probability, we may
therefore conclude that
\[
\norm{Z_n}_{L^1}\longrightarrow0.
\]

Now apply Lemma~\ref{lem:l1lower} to each row of the array.  The assumptions of
that lemma are satisfied because every $\xi_{n,i}$ is centered,
$L^2$-normalized, and obeys $\norm{\xi_{n,i}}_{L^1}\ge\delta$.  Hence
\[
\norm{Z_n}_{L^1}
=\Bigl\|\sum_{i\ge2} c_i^{(n)}\xi_{n,i}\Bigr\|_{L^1}
\ge
\frac{\delta}{2\sqrt2}
\Bigl(\sum_{i\ge2}(c_i^{(n)})^2\Bigr)^{1/2}
=\frac{\delta}{2\sqrt2}\norm{c^{(n)}}_2.
\]
The left-hand side tends to zero, so the right-hand side must also tend to
zero.  Therefore
\[
\norm{c^{(n)}}_2\longrightarrow0,
\]
which is the desired conclusion.
\end{proof}

\section{First Proof: Compactness}\label{sec:compactness-proof}

This section proves Theorem~\ref{thm:main}.  We follow the compactness route
from the manuscript `compactness\_formal.tex', but reorganize it so that the
head--tail structure and the endpoint alternatives are explicit.

\subsection{Contradiction array and profile decomposition}

\begin{proposition}[Normalized counterexample array]\label{prop:counterexample}
If Theorem~\ref{thm:main} fails, then there exist
\begin{enumerate}[label=\rm(\roman*)]
\item a triangular array $(\xi_{n,i})_{n\ge1,\ i\ge2}$ such that for each
fixed $n$ the family $(\eta,\xi_{n,i})_{i\ge2}$ is independent and every
$\xi_{n,i}$ satisfies
\[
\norm{\xi_{n,i}}_{L^2}=1,
\qquad
\E[\xi_{n,i}]=0,
\qquad
\delta\le \norm{\xi_{n,i}}_{L^1}\le 1-\delta;
\]
\item coefficient vectors $a^{(n)},b^{(n)}\in \ell^2(\N)$ with
\[
\norm{a^{(n)}}_2=\norm{b^{(n)}}_2=1,
\qquad
\ip{a^{(n)}}{b^{(n)}}=0,
\]
such that, writing
\[
X_n:=a_1^{(n)}\eta+\sum_{i\ge2} a_i^{(n)}\xi_{n,i},
\qquad
Y_n:=b_1^{(n)}\eta+\sum_{i\ge2} b_i^{(n)}\xi_{n,i},
\]
one has
\[
\norm{\abs{X_n}-\abs{Y_n}}_{L^2}\longrightarrow0.
\]
\end{enumerate}
\end{proposition}

\begin{proof}
Assume that the conclusion of Theorem~\ref{thm:main} is false.  Then there is
no constant $C$ for which the stable phase retrieval estimate holds uniformly
on the closed span of the coordinates.  By the orthogonal reduction
(Proposition~\ref{prop:orth-reduction}), this means that we may find, for each
$n\in\N$, an orthogonal pair of unit vectors in the span whose modulus gap is
at most $1/n$.  Concretely, there exist functions $X_n,Y_n$ in
$\overline{\Span}(\eta,\xi_i:i\ge2)$ such that
\[
\norm{X_n}_{L^2}=\norm{Y_n}_{L^2}=1,
\qquad
\E[X_nY_n]=0,
\qquad
\norm{|X_n|-|Y_n|}_{L^2}\le \frac1n.
\]

Because the family $(\eta,\xi_i)_{i\ge2}$ is orthonormal in $L^2$, we may write
\[
X_n=a_1^{(n)}\eta+\sum_{i\ge2}a_i^{(n)}\xi_i,
\qquad
Y_n=b_1^{(n)}\eta+\sum_{i\ge2}b_i^{(n)}\xi_i,
\]
with coefficient vectors $a^{(n)},b^{(n)}\in\ell^2(\N)$.  Parseval's identity
then yields
\[
\norm{a^{(n)}}_2=\norm{X_n}_{L^2}=1,
\qquad
\norm{b^{(n)}}_2=\norm{Y_n}_{L^2}=1,
\]
and, since the two functions are orthogonal in $L^2$,
\[
\ip{a^{(n)}}{b^{(n)}}=\E[X_nY_n]=0.
\]

Finally, to prepare for later rowwise permutations of the good coordinates, we
rename the good variables used in the $n$-th row by writing them as
$(\xi_{n,i})_{i\ge2}$.  This does not alter their laws or any of the stated
hypotheses, but it is convenient for the compactness scheme.  The displayed
modulus convergence is exactly our chosen violating sequence.  This proves the
proposition.
\end{proof}

\begin{lemma}[Reordered profile decomposition]\label{lem:profile}
After passing to a subsequence and, row by row, applying the same permutation
to the good coordinates and to their coefficients, one may find vectors
$a,b\in\ell^2(\N)$ and integers $m_n\to\infty$ such that, with
\[
a_{\mathrm{head}}^{(n)}:=a^{(n)}\mathbf 1_{\{1,\dots,m_n\}},
\qquad
a_{\mathrm{tail}}^{(n)}:=a^{(n)}\mathbf 1_{\{m_n+1,m_n+2,\dots\}},
\]
and likewise for $b^{(n)}$, the following hold:
\begin{enumerate}[label=\rm(\arabic*)]
\item $a_{\mathrm{head}}^{(n)}\to a$ and $b_{\mathrm{head}}^{(n)}\to b$
strongly in $\ell^2$;
\item
\[
\max_{i>m_n}\max(\abs{a_i^{(n)}},\abs{b_i^{(n)}})\to0;
\]
\item
\[
\norm{a_{\mathrm{tail}}^{(n)}}_2^2\to1-\norm{a}_2^2,
\qquad
\norm{b_{\mathrm{tail}}^{(n)}}_2^2\to1-\norm{b}_2^2,
\]
and
\[
\ip{a_{\mathrm{tail}}^{(n)}}{b_{\mathrm{tail}}^{(n)}}
\to -\ip{a}{b}.
\]
\end{enumerate}
\end{lemma}

\begin{proof}
For each row $n$ and each good index $i\ge2$, define
\[
s_i^{(n)}:=\max\bigl(\abs{a_i^{(n)}},\abs{b_i^{(n)}}\bigr).
\]
Permuting the good coordinates within the $n$-th row, we may assume that the
sequence $(s_i^{(n)})_{i\ge2}$ is nonincreasing in $i$.  The same permutation
is applied simultaneously to the variables $\xi_{n,i}$ and to both coefficient
vectors, so no hypothesis is changed.

The monotone rearrangement gives a uniform pointwise bound on the reordered
coefficients.  Indeed,
\[
\sum_{i\ge2}(s_i^{(n)})^2
\le
\sum_{i\ge2}\bigl((a_i^{(n)})^2+(b_i^{(n)})^2\bigr)
\le 2.
\]
Since $(s_i^{(n)})_{i\ge2}$ is nonincreasing, for every $i\ge2$ we have
\[
(i-1)(s_i^{(n)})^2
\le \sum_{j=2}^{i}(s_j^{(n)})^2
\le 2,
\]
and therefore
\[
s_i^{(n)}\le \sqrt{\frac{2}{i-1}}.
\]
In particular, for each fixed coordinate $i$, the sequences
$\bigl(a_i^{(n)}\bigr)_{n\ge1}$ and $\bigl(b_i^{(n)}\bigr)_{n\ge1}$ are
bounded.

We now extract coordinatewise limits.  First pass to a subsequence so that
$a_1^{(n)}$ and $b_1^{(n)}$ converge.  Next pass to a further subsequence so
that $a_2^{(n)}$ and $b_2^{(n)}$ converge, and continue inductively.  A
standard diagonal argument yields a single subsequence, not relabeled, for
which every fixed coordinate converges:
\[
a_i^{(n)}\to a_i,
\qquad
b_i^{(n)}\to b_i
\qquad(i\ge1).
\]
Set
\[
a:=(a_i)_{i\ge1},
\qquad
b:=(b_i)_{i\ge1}.
\]

We claim that $a,b\in\ell^2(\N)$.  Fix $M\in\N$.  Since the first $M$
coordinates converge,
\[
\sum_{i=1}^{M} a_i^2
=\lim_{n\to\infty}\sum_{i=1}^{M}(a_i^{(n)})^2
\le \liminf_{n\to\infty}\sum_{i\ge1}(a_i^{(n)})^2
=1.
\]
Letting $M\to\infty$ and using monotone convergence gives
$\sum_i a_i^2\le1$.  The same argument yields $\sum_i b_i^2\le1$.

Because $a,b\in\ell^2$, we may choose an increasing sequence of integers
$(m_k)_{k\ge1}$ with $m_k\to\infty$ such that
\[
\sum_{i>m_k} a_i^2+\sum_{i>m_k} b_i^2\le 2^{-k}
\qquad(k\ge1).
\]
For each fixed $k$, coordinatewise convergence implies
\[
\sum_{i=1}^{m_k}\abs{a_i^{(n)}-a_i}^2
+
\sum_{i=1}^{m_k}\abs{b_i^{(n)}-b_i}^2
\longrightarrow0
\qquad(n\to\infty).
\]
Passing to another subsequence if necessary, we may assume that for each $n$,
\[
\sum_{i=1}^{m_n}\abs{a_i^{(n)}-a_i}^2
+
\sum_{i=1}^{m_n}\abs{b_i^{(n)}-b_i}^2
\le 2^{-n}.
\]

We now verify the three conclusions.  First,
\[
\norm{a_{\mathrm{head}}^{(n)}-a}_2^2
\le
\sum_{i=1}^{m_n}\abs{a_i^{(n)}-a_i}^2
+
\sum_{i>m_n}a_i^2
\le 2^{-n}+2^{-n},
\]
so $a_{\mathrm{head}}^{(n)}\to a$ in $\ell^2$.  The same estimate gives
$b_{\mathrm{head}}^{(n)}\to b$.

Second, if $i>m_n$, then by monotonicity,
\[
\max\bigl(\abs{a_i^{(n)}},\abs{b_i^{(n)}}\bigr)
=s_i^{(n)}
\le s_{m_n+1}^{(n)}
\le \sqrt{\frac{2}{m_n}}.
\]
Since $m_n\to\infty$, this proves
\[
\max_{i>m_n}\max\bigl(\abs{a_i^{(n)}},\abs{b_i^{(n)}}\bigr)\to0.
\]

Finally, because each full coefficient vector has $\ell^2$ norm one,
\[
\norm{a_{\mathrm{tail}}^{(n)}}_2^2
=1-\norm{a_{\mathrm{head}}^{(n)}}_2^2
\longrightarrow 1-\norm{a}_2^2,
\]
and similarly
\[
\norm{b_{\mathrm{tail}}^{(n)}}_2^2
\longrightarrow 1-\norm{b}_2^2.
\]
For the inner products, use the orthogonality of the full vectors:
\[
0
=\ip{a^{(n)}}{b^{(n)}}
=\ip{a_{\mathrm{head}}^{(n)}}{b_{\mathrm{head}}^{(n)}}
+\ip{a_{\mathrm{tail}}^{(n)}}{b_{\mathrm{tail}}^{(n)}}.
\]
Since the head pieces converge strongly,
\[
\ip{a_{\mathrm{head}}^{(n)}}{b_{\mathrm{head}}^{(n)}}
\to \ip{a}{b},
\]
and therefore
\[
\ip{a_{\mathrm{tail}}^{(n)}}{b_{\mathrm{tail}}^{(n)}}
\to -\ip{a}{b}.
\]
This proves the lemma.
\end{proof}

\subsection{Head identification}

\begin{proposition}[Exceptional probes]\label{prop:exceptional-probes}
There exists a bounded measurable function $S_\eta$ such that
\[
\E[S_\eta]=0,
\qquad
\E[\eta S_\eta]=1.
\]
If $\eta^2\not\equiv1$ almost surely, then there exists a bounded measurable
function $T_\eta$ such that
\[
\E[T_\eta]=0,
\qquad
\E[(\eta^2-1)T_\eta]=1.
\]
\end{proposition}

\begin{proof}
Since $\norm{\eta}_{L^2}=1$, the random variable $\eta$ is not almost surely
zero.  Hence $\E[\abs{\eta}]>0$.  Applying the linear-probe construction from
Lemma~\ref{lem:probes} with $A=\E[\abs{\eta}]$ produces a bounded measurable
function $S_\eta$ satisfying
\[
\E[S_\eta]=0,
\qquad
\E[\eta S_\eta]=1.
\]

Assume now that $\eta^2\not\equiv1$ almost surely.  Then the centered random
variable $\eta^2-1$ is not almost surely zero, because
\[
\E[\eta^2-1]=\E[\eta^2]-1=0.
\]
Therefore its $L^1$ norm is strictly positive:
\[
\E[\abs{\eta^2-1}]>0.
\]
Applying again the linear-probe construction, now to the random variable
$\eta^2-1$, we obtain a bounded measurable function $T_\eta$ such that
\[
\E[T_\eta]=0,
\qquad
\E[(\eta^2-1)T_\eta]=1.
\]
This is exactly the desired conclusion.
\end{proof}

\begin{lemma}[Identification of the head profile]\label{lem:head-sign}
Either there exists $\tau\in\{\pm1\}$ such that
\[
a=\tau b
\qquad\text{and}\qquad
\norm{a_{\mathrm{head}}^{(n)}-\tau b_{\mathrm{head}}^{(n)}}_2\to0,
\]
or else $\eta$ is a Rademacher variable and
\[
a_i=b_i=0
\qquad(i\ge2).
\]
\end{lemma}

\begin{proof}
We first convert the convergence of moduli into convergence of the quadratic
observable.  Since
\[
\norm{X_n}_{L^2}=\norm{Y_n}_{L^2}=1,
\]
we have
\[
\norm{|X_n|+|Y_n|}_{L^2}\le 2.
\]
Therefore
\[
\norm{|X_n|^2-|Y_n|^2}_{L^1}
=\norm{(|X_n|-|Y_n|)(|X_n|+|Y_n|)}_{L^1}
\le 2\norm{|X_n|-|Y_n|}_{L^2}
\longrightarrow0.
\]
Because the variables are real-valued, $|X_n|^2=X_n^2$ and $|Y_n|^2=Y_n^2$,
so we have established
\[
\norm{X_n^2-Y_n^2}_{L^1}\longrightarrow0.
\]

\smallskip
\noindent\emph{Step 1: diagonal good-coordinate identities.}
Fix $i\ge2$ and let $T_{n,i}$ be the quadratic probe associated with the
coordinate $\xi_{n,i}$.  The probes are uniformly bounded in $L^\infty$, so the
$L^1$ convergence above implies
\[
\E[(X_n^2-Y_n^2)T_{n,i}]\longrightarrow0.
\]
We compute the left-hand side explicitly.  Expanding $X_n^2$ and integrating
against $T_{n,i}$, every term involving a coordinate different from
$\xi_{n,i}$ vanishes because $T_{n,i}$ depends only on $\xi_{n,i}$ and is
centered.  Mixed terms involving $\xi_{n,i}$ and another coordinate also
vanish by independence and centering.  The only surviving contribution is the
term $(a_i^{(n)})^2\xi_{n,i}^2$.  Hence
\[
\E[X_n^2T_{n,i}]
=(a_i^{(n)})^2\E[\xi_{n,i}^2T_{n,i}]
=(a_i^{(n)})^2,
\]
because
\[
\E[\xi_{n,i}^2T_{n,i}]
=\E[(\xi_{n,i}^2-1)T_{n,i}]+\E[T_{n,i}]
=1+0.
\]
Similarly,
\[
\E[Y_n^2T_{n,i}]=(b_i^{(n)})^2.
\]
Therefore
\[
(a_i^{(n)})^2-(b_i^{(n)})^2
=\E[(X_n^2-Y_n^2)T_{n,i}]
\longrightarrow0.
\]
Passing to the limit in $n$ gives
\[
a_i^2=b_i^2
\qquad(i\ge2).
\]

\smallskip
\noindent\emph{Step 2: off-diagonal good-good identities.}
Fix distinct indices $i,j\ge2$ and let $S_{n,i},S_{n,j}$ be the corresponding
linear probes.  Again boundedness and $L^1$ convergence imply
\[
\E[(X_n^2-Y_n^2)S_{n,i}S_{n,j}]\longrightarrow0.
\]
In the expansion of $X_n^2$, all terms vanish after integration against
$S_{n,i}S_{n,j}$ except the mixed term
$2a_i^{(n)}a_j^{(n)}\xi_{n,i}\xi_{n,j}$.  Indeed, any term missing one of the
coordinates $\xi_{n,i}$ or $\xi_{n,j}$ is killed by the zero mean of the
corresponding probe, and terms involving other coordinates factor through zero
means by independence.  Hence
\[
\E[X_n^2S_{n,i}S_{n,j}]
=2a_i^{(n)}a_j^{(n)}
\E[\xi_{n,i}S_{n,i}]\,\E[\xi_{n,j}S_{n,j}]
=2a_i^{(n)}a_j^{(n)}.
\]
Similarly,
\[
\E[Y_n^2S_{n,i}S_{n,j}]
=2b_i^{(n)}b_j^{(n)}.
\]
Thus
\[
2(a_i^{(n)}a_j^{(n)}-b_i^{(n)}b_j^{(n)})\longrightarrow0,
\]
and passing to the limit gives
\[
a_ia_j=b_ib_j
\qquad(i\ne j,\ i,j\ge2).
\]

\smallskip
\noindent\emph{Step 3: exceptional-good identities.}
Fix $j\ge2$.  Let $S_\eta$ be the bounded probe from
Proposition~\ref{prop:exceptional-probes}.  By the same reasoning,
\[
\E[(X_n^2-Y_n^2)S_\eta S_{n,j}]\longrightarrow0.
\]
Once more only one mixed term survives, namely
$2a_1^{(n)}a_j^{(n)}\eta\xi_{n,j}$ for $X_n^2$ and the corresponding term for
$Y_n^2$.  Therefore
\[
\E[X_n^2S_\eta S_{n,j}]
=2a_1^{(n)}a_j^{(n)}\E[\eta S_\eta]\E[\xi_{n,j}S_{n,j}]
=2a_1^{(n)}a_j^{(n)},
\]
and analogously
\[
\E[Y_n^2S_\eta S_{n,j}]
=2b_1^{(n)}b_j^{(n)}.
\]
Hence
\[
a_1a_j=b_1b_j
\qquad(j\ge2).
\]

\smallskip
\noindent\emph{Step 4: deduction of the sign.}
Assume first that there exists some index $i_0\ge2$ such that $a_{i_0}\ne0$.
From the diagonal identity we know that
$a_{i_0}^2=b_{i_0}^2$, so $b_{i_0}=\tau a_{i_0}$ for some
$\tau\in\{\pm1\}$.  Then for every $j\ge2$, the off-diagonal identity gives
\[
a_{i_0}a_j=b_{i_0}b_j=\tau a_{i_0}b_j.
\]
Since $a_{i_0}\ne0$, we conclude that $a_j=\tau b_j$ for every $j\ge2$.
Applying the exceptional-good identity with $j=i_0$, we obtain
\[
a_1a_{i_0}=b_1b_{i_0}=\tau a_{i_0}b_1,
\]
and again $a_{i_0}\ne0$ implies $a_1=\tau b_1$.  Hence $a=\tau b$.

Assume next that
\[
a_i=0
\qquad\text{for every }i\ge2.
\]
Then the diagonal identities show that also $b_i=0$ for every $i\ge2$.
If $\eta^2\not\equiv1$ almost surely, let $T_\eta$ be the quadratic probe from
Proposition~\ref{prop:exceptional-probes}.  Repeating the diagonal argument for
the exceptional coordinate gives
\[
(a_1^{(n)})^2-(b_1^{(n)})^2
=\E[(X_n^2-Y_n^2)T_\eta]
\longrightarrow0,
\]
whence $a_1^2=b_1^2$ and therefore $a=\tau b$ for some sign $\tau$.

The only remaining possibility is that $\eta^2\equiv1$ almost surely.  Since
$\E[\eta]=0$ and $\norm{\eta}_{L^2}=1$, this means exactly that $\eta$ is a
Rademacher variable.  In that case we are in the second alternative stated in
the lemma.

Finally, if the first alternative holds, then the convergence of the head
profiles follows from Lemma~\ref{lem:profile}:
\[
\norm{a_{\mathrm{head}}^{(n)}-\tau b_{\mathrm{head}}^{(n)}}_2
\le
\norm{a_{\mathrm{head}}^{(n)}-a}_2
+
\norm{b_{\mathrm{head}}^{(n)}-b}_2
\longrightarrow0.
\]
\end{proof}

\subsection{Limit law and exact modulus identity}

\begin{lemma}[Limit law]\label{lem:limit-law}
After passing to a subsequence, the quadruple
\[
\bigl(X_n^{\mathrm{head}},Y_n^{\mathrm{head}},
X_n^{\mathrm{tail}},Y_n^{\mathrm{tail}}\bigr)
\]
converges in distribution to a random vector
\[
(H_X,H_Y,R_X,R_Y)\in\R^4
\]
such that:
\begin{enumerate}[label=\rm(\arabic*)]
\item $(H_X,H_Y)$ is independent of $(R_X,R_Y)$;
\item the law of $(R_X,R_Y)$ is infinitely divisible;
\item either $H_X=\tau H_Y$ almost surely for some $\tau\in\{\pm1\}$, or
else $\eta$ is Rademacher and
$H_X=a_1\eta$, $H_Y=b_1\eta$ almost surely;
\item
\[
\abs{H_X+R_X}=\abs{H_Y+R_Y}
\qquad\text{almost surely.}
\]
\end{enumerate}
\end{lemma}

\begin{proof}
The head variables converge in $L^2$ by the strong convergence of the head
coefficients and the orthonormality of each row.  The tail coefficient arrays
are infinitesimal by Lemma~\ref{lem:profile}, so every weak limit of the
tail pair is infinitely divisible.  Independence of head and tail persists
because they are built from disjoint independent coordinates.  Finally,
$\norm{\abs{X_n}-\abs{Y_n}}_{L^2}\to0$ implies
$\abs{X_n}-\abs{Y_n}\to0$ in probability; passing to the joint weak limit
gives the exact identity in the limit.  The description of the head follows
from Lemma~\ref{lem:head-sign}.
\end{proof}

\begin{lemma}[Infinitely divisible laws on the phase-retrieval cone are
line-supported]\label{lem:id-line}
Let $\mu$ be an infinitely divisible probability measure on $\R^2$ with
\[
\supp\mu\subset\Gamma:=\{(u,v)\in\R^2:\abs{u}=\abs{v}\}.
\]
Then there exists $\tau\in\{\pm1\}$ such that
\[
\supp\mu\subset L_\tau:=\{(u,\tau u):u\in\R\}.
\]
\end{lemma}

\begin{proof}
The set $\Gamma$ is exactly the union $L_+\cup L_-$.  Since $\mu$ is
infinitely divisible, it is in particular $2$-divisible: there exists a
probability measure $\nu$ such that $\mu=\nu*\nu$.  Let $A=\supp\nu$.  Then
\[
\supp\mu=\overline{A+A}\subset L_+\cup L_-.
\]
If $x,y\in A$ are distinct, then $2x,x+y,2y\in A+A$.  If $2x$ and $2y$ belong
to the same line $L_\tau$, then so does $x+y$, forcing $x,y\in L_\tau$.  If
they belong to different lines, then the midpoint $x+y$ is the origin, so
$y=-x$ and again $x,y$ lie in one of the lines.  Thus every two points of $A$
lie on a common phase-retrieval line.  Since $L_+\cap L_-=\{0\}$, this implies
that $A$ is contained in one of the two lines, and so is $\supp\mu$.
\end{proof}

\subsection{End of the compactness proof}

\begin{proof}[Proof of Theorem~\ref{thm:main}]
Assume that Theorem~\ref{thm:main} fails and let the counterexample array be as
in Proposition~\ref{prop:counterexample}.  Apply
Lemmas~\ref{lem:profile}, \ref{lem:head-sign}, and \ref{lem:limit-law}.

If the first alternative of Lemma~\ref{lem:head-sign} holds, fix
$\tau\in\{\pm1\}$ such that $H_Y=\tau H_X$ almost surely.  The exact modulus
identity becomes
\[
\abs{H_X+R_X}=\abs{\tau H_X+R_Y}
\qquad\text{almost surely.}
\]
If $H_X$ is not almost surely zero, then conditioning on the tail and varying
the head over two support points of $H_X$ shows that
$R_Y-\tau R_X=0$ almost surely.  If $H_X\equiv0$, then
$(R_X,R_Y)$ itself satisfies $\abs{R_X}=\abs{R_Y}$ almost surely, so
Lemma~\ref{lem:id-line} yields
$R_Y=\tau_0 R_X$ almost surely for some $\tau_0\in\{\pm1\}$.

In either case, for a suitable sign $\sigma\in\{\pm1\}$ the random variables
\[
Z_n:=\sum_{i>m_n}(b_i^{(n)}-\sigma a_i^{(n)})\xi_{n,i}
\]
converge in distribution to $0$.  Since the coefficients are diffuse,
Lemma~\ref{lem:tail-rigidity} gives
\[
\norm{b_{\mathrm{tail}}^{(n)}-\sigma a_{\mathrm{tail}}^{(n)}}_2\to0.
\]
Together with Lemma~\ref{lem:head-sign} this implies
$\norm{b^{(n)}-\sigma a^{(n)}}_2\to0$, contradicting
\[
\norm{b^{(n)}-\sigma a^{(n)}}_2^2
=\norm{a^{(n)}}_2^2+\norm{b^{(n)}}_2^2
-2\sigma\ip{a^{(n)}}{b^{(n)}}=2.
\]

If the second alternative of Lemma~\ref{lem:head-sign} holds, then
$\eta$ is Rademacher and the head limits are multiples of $\eta$ alone.  The
same argument applies, since in that case the good coordinates disappear from
the head and the tail limit still lies on $\Gamma$.
\end{proof}

\subsection{A Gaussian alternative in the degenerate-head regime}

The previous endpoint used only the support structure of infinitely divisible
limits.  The next result records a second route, adapted from the complex
one-exception manuscript, showing that in the degenerate-head regime the
diffuse tail automatically generates a nontrivial Gaussian component.

\begin{proposition}[Diffuse tails force a nontrivial Gaussian part]
\label{prop:gaussian-part}
Let
\[
V_{n,i}:=(a_i^{(n)}\xi_{n,i},\,b_i^{(n)}\xi_{n,i})\in\R^2
\qquad(i>m_n),
\]
where the coordinates are independent, centered, $L^2$-normalized, and satisfy
$\norm{\xi_{n,i}}_{L^1}\ge \delta$.  Set
\[
\alpha_{n,i}:=\bigl((a_i^{(n)})^2+(b_i^{(n)})^2\bigr)^{1/2}.
\]
Assume
\[
\max_{i>m_n}\alpha_{n,i}\to0,
\qquad
\sum_{i>m_n}\alpha_{n,i}^2\to \sigma^2>0,
\]
and that the row sums $\sum_{i>m_n}V_{n,i}$ converge in distribution to an
infinitely divisible law $\mu$ on $\R^2$.  Then the Gaussian part of $\mu$ is
nontrivial.
\end{proposition}

\begin{proof}
Set $\beta_n=\max_{i>m_n}\alpha_{n,i}$ and $r_n=\beta_n^{1/2}$.  For each
$i>m_n$ with $\alpha_{n,i}\ne0$, let $K_{n,i}=r_n/\alpha_{n,i}\to\infty$
uniformly and define the centered truncation
\[
\widetilde\xi_{n,i}
:=
\xi_{n,i}\mathbf 1_{\{\abs{\xi_{n,i}}\le K_{n,i}\}}
-\E[\xi_{n,i}\mathbf 1_{\{\abs{\xi_{n,i}}\le K_{n,i}\}}].
\]
Since $\E[\abs{\xi_{n,i}}]\ge\delta$ and $\E[\xi_{n,i}^2]=1$, the same estimate
as in the complex proof shows that there exists $c_0=c_0(\delta)>0$ such that
\[
\E[\widetilde\xi_{n,i}^2]\ge c_0
\]
for all large $n$ and all $i>m_n$.  Define
$\widetilde V_{n,i}=(a_i^{(n)}\widetilde\xi_{n,i},\,b_i^{(n)}\widetilde\xi_{n,i})$.
Then $\abs{\widetilde V_{n,i}}\le 2r_n$, so the truncated array is bounded and
infinitesimal, and
\[
\sum_{i>m_n}\E[\abs{\widetilde V_{n,i}}^2]
\ge c_0\sum_{i>m_n}\alpha_{n,i}^2
\longrightarrow c_0\sigma^2>0.
\]
For infinitesimal triangular arrays, the covariance matrices of shrinking
centered truncations converge to the Gaussian covariance of the limit law; see
\cite[Chapter~IV]{gnedenko-kolmogorov}.  Since the traces of the truncation
covariances stay bounded away from zero, the Gaussian covariance of $\mu$ is
nonzero.
\end{proof}

\begin{remark}[Alternative Gaussian endpoint]\label{rem:gaussian-endpoint}
Assume we are in the degenerate-head regime $H_X\equiv H_Y\equiv0$, so the
tail law $\mu$ of $(R_X,R_Y)$ is supported on $\Gamma=\{(u,v):\abs{u}=\abs{v}\}$.
Write $\mu=\gamma*\nu$, where $\gamma$ is the Gaussian part and $\nu$ the
independent remainder.  By Proposition~\ref{prop:gaussian-part}, $\gamma$ is
nontrivial.  Its support is a nonzero linear subspace $E\subset\R^2$.

Because $\supp\mu=\supp\gamma+\supp\nu\subset\Gamma=L_+\cup L_-$, the
connected set $E$ must itself lie in one of the lines $L_\tau$.  Fix such a
$\tau$.  For any $z\in\supp\nu$, the translated line $E+z$ is contained in
$\Gamma$.  Since a translate of $L_\tau$ not equal to $L_\tau$ intersects
$L_+\cup L_-$ in at most two points, one must have $z\in L_\tau$.  Hence
$\supp\nu\subset L_\tau$ and therefore $\supp\mu\subset L_\tau$.  Thus
$R_Y=\tau R_X$ almost surely.  This recovers the final contradiction without
appealing to Lemma~\ref{lem:id-line}.
\end{remark}

\section{Second Proof: Anticoncentration}\label{sec:quant-proof}

This section gives a second proof of Theorem~\ref{thm:main}.  The argument is
quantitative in spirit.  Instead of passing to a limit law, we split the
coefficients into a head and a tail and show directly that one of two explicit
mechanisms forces a positive modulus gap.

By Proposition~\ref{prop:orth-reduction}, it is enough to prove an orthogonal
lower bound for arbitrary unit vectors in the one-exception span.  Thus let
$a=(a_i)_{i\ge1}, b=(b_i)_{i\ge1}\in \ell^2(\N)$ satisfy
\[
\norm{a}_2=\norm{b}_2=1,
\qquad
\ip{a}{b}=0,
\]
and set
\[
X:=a_1\eta+\sum_{i\ge2} a_i\xi_i,
\qquad
Y:=b_1\eta+\sum_{i\ge2} b_i\xi_i.
\]
Because the family $(\eta,\xi_i)_{i\ge2}$ is orthonormal in $L^2$, these sums
converge in $L^2$.  We shall prove that
\[
\norm{|X|-|Y|}_{L^2}\ge c_{\delta,\eta}>0,
\]
where the constant depends only on $\delta$ and on the law of $\eta$.

\subsection{Setup and the dichotomy}

Define
\[
\eps:=\E[\abs{X^2-Y^2}],
\qquad
\Delta:=\norm{|X|-|Y|}_{L^2}.
\]
As before,
\[
\eps\le 2\Delta.
\]

Fix
\[
\theta:=\frac{\delta^2}{8},
\qquad
\lambda:=\frac{\delta^3}{2^{20}},
\qquad
s_0:=\frac{\delta^4}{2^{20}}.
\]
Let $S_\eta$ be the bounded linear probe from
Proposition~\ref{prop:exceptional-probes}.  If $\eta^2\not\equiv1$, let
$T_\eta$ be the bounded quadratic probe from the same proposition.  Choose a
constant $K_{\delta,\eta}\ge1$ so large that
\[
K_{\delta,\eta}\ge \frac{4}{\delta^2},
\qquad
K_{\delta,\eta}\ge \frac{2}{\delta},
\qquad
K_{\delta,\eta}\ge \norm{S_\eta}_{L^\infty},
\]
and, if $\eta^2\not\equiv1$,
\[
K_{\delta,\eta}\ge \norm{T_\eta}_{L^\infty}.
\]

Define the head set by
\[
H=H(a,b):=\{1\}\cup\{i\ge2:\max(|a_i|,|b_i|)>\lambda\}.
\]
Then
\[
\#H\le d^*(\delta):=3+2^{41}\delta^{-6},
\]
and every tail coordinate $i\notin H$ satisfies
\[
\max(|a_i|,|b_i|)\le \lambda.
\]

The first step is to separate the proof into two mutually exclusive regimes.
Either one of the two signs makes the tail nearly one-dimensional, in which
case the argument is driven by head rigidity, or both signs retain substantial
tail mass, in which case anticoncentration becomes effective.

\begin{lemma}[Tail sign dichotomy]\label{lem:dichotomy}
Exactly one of the following alternatives holds:
\begin{enumerate}[label=\rm(\roman*)]
\item \textbf{Diffuse:}
\[
\norm{(a-b)\mathbf 1_{H^c}}_2^2\ge16s_0
\quad\text{and}\quad
\norm{(a+b)\mathbf 1_{H^c}}_2^2\ge16s_0.
\]
\item \textbf{Concentrated:} there exists $\tau\in\{\pm1\}$ such that
\[
\norm{(a-\tau b)\mathbf 1_{H^c}}_2<\theta/32.
\]
\end{enumerate}
\end{lemma}

\begin{proof}
We first note that
\[
16s_0=\Bigl(\frac{\theta}{32}\Bigr)^2.
\]
Thus the concentrated alternative may be rewritten as saying that there exists
$\tau\in\{\pm1\}$ such that
\[
\norm{(a-\tau b)\mathbf 1_{H^c}}_2^2<16s_0.
\]

If this happens for $\tau=1$ or $\tau=-1$, then we are in the concentrated
case and there is nothing more to prove.  Suppose, on the contrary, that the
concentrated alternative fails.  Then for both choices of sign we have
\[
\norm{(a-b)\mathbf 1_{H^c}}_2^2\ge16s_0
\qquad\text{and}\qquad
\norm{(a+b)\mathbf 1_{H^c}}_2^2\ge16s_0.
\]
But this is exactly the diffuse alternative.

Finally, the two alternatives are mutually exclusive, because if the
concentrated alternative held for some sign $\tau$, then one of the two vectors
$a-b$ or $a+b$ would have tail norm squared strictly smaller than $16s_0$,
which contradicts the diffuse inequalities.  Hence exactly one of the two
alternatives holds.
\end{proof}

\subsection{The concentrated branch}

We now turn to the first regime.  In the concentrated case, the tail of
$a-\tau b$ is small for a suitable sign $\tau$, so any substantial discrepancy
between the two vectors must be visible on the head.  The next proposition
shows that the quadratic observable $X^2-Y^2$ controls precisely those
head coefficients.

\begin{proposition}[Head entry bounds with one exceptional coordinate]
\label{prop:head-entry}
The following estimates hold:
\begin{enumerate}[label=\rm(\arabic*)]
\item if $i,j\in H\cap\{2,3,\dots\}$, then
\[
|a_i a_j-b_i b_j|\le K_{\delta,\eta}\,\eps;
\]
\item if $j\in H\cap\{2,3,\dots\}$, then
\[
|a_1a_j-b_1b_j|\le K_{\delta,\eta}\,\eps;
\]
\item if $\eta^2\not\equiv1$, then
\[
|a_1^2-b_1^2|\le K_{\delta,\eta}\,\eps.
\]
\end{enumerate}
\end{proposition}

\begin{proof}
For $i=j\ge2$, use the quadratic probe $T_i$ for $\xi_i$:
\[
|a_i^2-b_i^2|
=\bigl|\E[(X^2-Y^2)T_i]\bigr|
\le \norm{T_i}_{L^\infty}\eps
\le K_{\delta,\eta}\eps.
\]
For distinct $i,j\ge2$, use the linear probes $S_i,S_j$:
\[
2|a_i a_j-b_i b_j|
=\bigl|\E[(X^2-Y^2)S_iS_j]\bigr|
\le \norm{S_i}_{L^\infty}\norm{S_j}_{L^\infty}\eps
\le K_{\delta,\eta}\eps.
\]
This proves \rm(1).

For \rm(2), use the product $S_\eta S_j$.  Exactly as in the proof of
Lemma~\ref{lem:head-sign}, all terms vanish except the mixed
exceptional-good term, so
\[
2|a_1a_j-b_1b_j|
=\bigl|\E[(X^2-Y^2)S_\eta S_j]\bigr|
\le \norm{S_\eta}_{L^\infty}\norm{S_j}_{L^\infty}\eps
\le K_{\delta,\eta}\eps.
\]

Finally, if $\eta^2\not\equiv1$, use $T_\eta$:
\[
|a_1^2-b_1^2|
=\bigl|\E[(X^2-Y^2)T_\eta]\bigr|
\le \norm{T_\eta}_{L^\infty}\eps
\le K_{\delta,\eta}\eps.
\]
\end{proof}

To convert these entrywise bounds into a global contradiction, we use the same
rank-one algebra that appeared in the homogeneous proof.  This isolates the
geometric fact that two orthogonal unit vectors cannot agree too well on the
head unless the corresponding head difference matrix is small.

\begin{lemma}[Rank-one algebra]\label{lem:rank-one}
For $x,y\in\R^m$, let $D(x,y)_{ij}=x_ix_j-y_iy_j$.  Then
\[
\norm{x-y}_2^2\,\norm{x+y}_2^2\le 2\norm{D(x,y)}_F^2.
\]
\end{lemma}

\begin{proof}
Writing $s=\norm{x}_2^2$, $t=\norm{y}_2^2$, and $c=\ip{x}{y}$, the left-hand
side equals $(s+t)^2-4c^2$, while the right-hand side is
$2(s^2+t^2-2c^2)$.  Their difference is $(s-t)^2$, which is nonnegative.
\end{proof}

Combining the head-entry bounds with the rank-one inequality yields the
concentrated branch.  The proof splits off the genuinely degenerate
Rademacher-only head, since only there can the exceptional coordinate fail to
provide a quadratic probe.

\begin{proposition}[Concentrated branch]\label{prop:concentrated}
If there exists $\tau\in\{\pm1\}$ such that
\[
\norm{(a-\tau b)\mathbf 1_{H^c}}_2\le \theta/32,
\]
then
\[
\eps\ge c_{\mathrm{line}}(\delta,\eta)>0.
\]
\end{proposition}

\begin{proof}
We distinguish two cases.

\smallskip
\noindent\emph{Case 1: $\eta$ is not Rademacher, or $H\cap\{2,3,\dots\}$ is
nonempty.}
In this case the head coefficients are rigidly linked by
Proposition~\ref{prop:head-entry}: either directly through the diagonal
estimate for the exceptional coordinate or, in the Rademacher case, through a
good head coordinate exactly as in the proof of Lemma~\ref{lem:head-sign}.  In
particular there exists a sign, which must be the same $\tau$, such that the
full head difference matrix of $a\mathbf 1_H$ and $\tau b\mathbf 1_H$ has all
entries bounded by $K_{\delta,\eta}\eps$.

Let
\[
x:=a\mathbf 1_H,
\qquad
y:=\tau b\mathbf 1_H.
\]
Applying Lemma~\ref{lem:rank-one},
\[
\norm{x-y}_2^2\,\norm{x+y}_2^2
\le 2\#H^2 K_{\delta,\eta}^2\eps^2.
\]
Since the tail of $a-\tau b$ has norm at most $\theta/32$ and
$\norm{a-\tau b}_2^2=2$ by orthogonality, the head part satisfies
\[
\norm{x-y}_2^2
=2-\norm{(a-\tau b)\mathbf 1_{H^c}}_2^2
\ge 1.
\]
Hence
\[
\norm{x+y}_2\le 2\#H K_{\delta,\eta}\eps
\le 2d^*(\delta)K_{\delta,\eta}\eps.
\]
If $\eps$ were smaller than a sufficiently small constant
$c_{\mathrm{line}}(\delta,\eta)$, then the head of $a+\tau b$ would be tiny
while the tail of $a-\tau b$ is already tiny.  The same gluing argument as in
the homogeneous proof then yields a contradiction.  Therefore
$\eps\ge c_{\mathrm{line}}(\delta,\eta)$.

\smallskip
\noindent\emph{Case 2: $\eta$ is Rademacher and $H=\{1\}$.}
Then all good coordinates lie in the tail and satisfy
$|a_i|,|b_i|\le \lambda$ for $i\ge2$.  The concentrated assumption gives a sign
$\tau$ such that $(a-\tau b)\mathbf 1_{H^c}$ is small.  If $\tau=1$, then the
good parts of $a$ and $b$ are almost equal, so
\[
\ip{a}{b}
=a_1b_1+\ip{a\mathbf 1_{H^c}}{b\mathbf 1_{H^c}}
\]
has the same sign as $a_1b_1$ up to a small error, contradicting
$\ip{a}{b}=0$ once the constants are fixed.  The case $\tau=-1$ is analogous.
Thus this degenerate Rademacher subcase cannot occur.
\end{proof}

\subsection{The diffuse branch}

We now move to the complementary regime.  Here both $a-b$ and $a+b$ retain a
definite amount of $\ell^2$ mass on the tail, and every tail coordinate is
individually small.  This is exactly the setting in which block
anticoncentration can be used.

Set
\[
c:=\frac{(a-b)\mathbf 1_{H^c}}{\sqrt2},
\qquad
d:=\frac{(a+b)\mathbf 1_{H^c}}{\sqrt2}.
\]
Because $1\in H$, the vectors $c$ and $d$ are supported only on good
coordinates.  In the diffuse case both have $\ell^2$ norm at least
$\sqrt{8s_0}$, while each coordinate is bounded by $\sqrt2\,\lambda$.

The next lemma packages the quantitative probabilistic input.  It says that
after arbitrary deterministic shifts, the two tail combinations cannot both
remain close to zero with high probability.

\begin{lemma}[Quantitative cross anticoncentration]\label{lem:cross}
There exists $r_*=r_*(\delta)>0$ such that, in the diffuse case,
\[
\Prob\bigl(r_*<\min(\abs{X_c+u},\abs{X_d+v})\bigr)\ge \frac14
\qquad(u,v\in\R).
\]
\end{lemma}

\begin{proof}
We prove first a one-variable small-ball estimate.  Let $e\in \ell^2(\N)$ be
supported on $H^c$ and satisfy
\[
\norm{e}_2^2\ge 8s_0,
\qquad
\sup_{i\ge2}\abs{e_i}\le \sqrt2\,\lambda.
\]
We claim that there exists $r_0=r_0(\delta)>0$ such that
\begin{equation}\label{eq:onevar-smallball}
\Prob\bigl(\abs{X_e-t}\le r_0\bigr)\le \frac38
\qquad(t\in\R).
\end{equation}
Once this is proved, the lemma follows immediately by applying
\eqref{eq:onevar-smallball} to $e=c$ with $t=-u$ and to $e=d$ with $t=-v$.
Indeed, in the diffuse case both $c$ and $d$ satisfy the displayed
hypotheses, so
\[
\Prob\bigl(\abs{X_c+u}\le r_0\bigr)\le \frac38,
\qquad
\Prob\bigl(\abs{X_d+v}\le r_0\bigr)\le \frac38.
\]
By the union bound,
\[
\Prob\bigl(\abs{X_c+u}\le r_0\ \text{or}\ \abs{X_d+v}\le r_0\bigr)
\le \frac34,
\]
and therefore
\[
\Prob\bigl(r_0<\min(\abs{X_c+u},\abs{X_d+v})\bigr)\ge \frac14.
\]
Thus it remains only to establish \eqref{eq:onevar-smallball}.

\smallskip
\noindent\emph{Step 1: a block partition.}
Set
\[
M:=\Bigl\lceil 2^{16}\delta^{-2}\Bigr\rceil,
\qquad
\eta:=\frac{s_0}{2M},
\qquad
r_0:=\frac{\delta\sqrt{\eta}}{16}.
\]
Because
\[
2\lambda^2=\frac{\delta^6}{2^{39}}
\le \frac{\delta^6}{2^{38}}
\le \eta,
\]
every coordinate satisfies $e_i^2\le \eta$.  Since
\[
\norm{e}_2^2\ge 8s_0>2M\eta=s_0,
\]
we may choose pairwise disjoint finite sets $G_1,\dots,G_M\subset H^c$ such
that
\begin{equation}\label{eq:block-energy}
\eta\le \sum_{i\in G_j} e_i^2\le 2\eta
\qquad(1\le j\le M).
\end{equation}
Indeed, construct the blocks greedily.  Suppose $G_1,\dots,G_{j-1}$ have
already been chosen.  Their total energy is at most $2(j-1)\eta$, so the
unused part of $e$ still has energy at least
\[
8s_0-2(j-1)\eta\ge 8s_0-2M\eta=7s_0> \eta.
\]
Hence one can choose a finite subset of the remaining support whose energy is
at least $\eta$, and by minimality its energy is at most $\eta+\sup_i e_i^2
\le 2\eta$.  This produces the next block.

\smallskip
\noindent\emph{Step 2: symmetrized block sums have uniform two-sided mass.}
Let $X_e'$ be an independent copy of $X_e$, realized on a second copy of the
probability space.  For each $i\ge2$ define
\[
D_i:=\xi_i(\omega_1)-\xi_i(\omega_2).
\]
Then the variables $(D_i)_{i\ge2}$ are independent, centered, and symmetric,
with
\[
\norm{D_i}_{L^2}^2
=\E[(\xi_i-\xi_i')^2]
=2.
\]
Moreover, by Jensen's inequality applied conditionally on $\xi_i$,
\[
\E[\abs{D_i}]
=\E\bigl[\E[\abs{\xi_i-\xi_i'}\mid \xi_i]\bigr]
\ge \E\bigl[\abs{\xi_i-\E[\xi_i']} \bigr]
=\E[\abs{\xi_i}]
\ge \delta.
\]
Hence
\[
\E[\abs{D_i}]
\ge \frac{\delta}{\sqrt2}\norm{D_i}_{L^2}.
\]
For each block define
\[
Y_j:=\sum_{i\in G_j} e_iD_i,
\qquad
\sigma_j^2:=\sum_{i\in G_j} e_i^2.
\]
By \eqref{eq:block-energy},
\[
\eta\le \sigma_j^2\le 2\eta.
\]
Applying Lemma~\ref{lem:l1lower} to the independent centered family
$(D_i)_{i\in G_j}$, with lower $L^1$ ratio $\delta/\sqrt2$, gives
\[
\E[\abs{Y_j}]
\ge
\frac{(\delta/\sqrt2)}{2\sqrt2}
\Bigl(\sum_{i\in G_j} e_i^2\norm{D_i}_{L^2}^2\Bigr)^{1/2}
=\frac{\delta}{4}\bigl(2\sigma_j^2\bigr)^{1/2}
=\frac{\delta}{2}\sigma_j
\ge \frac{\delta\sqrt{\eta}}{2}.
\]
Set
\[
a_0:=\frac{\delta\sqrt{\eta}}{4}=4r_0.
\]
Then $a_0\le \E[\abs{Y_j}]/2$.  Let
\[
A_j:=\{\abs{Y_j}\ge a_0\}.
\]
Since
\[
\E[Y_j^2]=2\sigma_j^2\le 4\eta,
\]
we obtain
\[
\E[\abs{Y_j}]
\le a_0+\E[\abs{Y_j}\mathbf 1_{A_j}]
\le a_0+\Prob(A_j)^{1/2}\E[Y_j^2]^{1/2}.
\]
Because $\E[\abs{Y_j}]\ge 2a_0$, the second term on the right is at least
$a_0$, so
\[
\Prob(A_j)^{1/2}\,(4\eta)^{1/2}\ge a_0.
\]
Therefore
\begin{equation}\label{eq:block-active}
\Prob(A_j)\ge \frac{a_0^2}{4\eta}=\frac{\delta^2}{64}.
\end{equation}

\smallskip
\noindent\emph{Step 3: many active blocks with high probability.}
Let
\[
N:=\sum_{j=1}^M \mathbf 1_{A_j}.
\]
The variables $\mathbf 1_{A_j}$ are independent because the blocks are
disjoint.  By \eqref{eq:block-active},
\[
\E[N]
=\sum_{j=1}^M \Prob(A_j)
\ge M\frac{\delta^2}{64}
\ge 1024.
\]
Since $\operatorname{Var}(\mathbf 1_{A_j})\le \E[\mathbf 1_{A_j}]$, we also have
\[
\operatorname{Var}(N)\le \E[N].
\]
Chebyshev's inequality gives
\[
\Prob(N<256)
\le
\Prob\bigl(\abs{N-\E[N]}\ge \E[N]-256\bigr)
\le
\frac{\operatorname{Var}(N)}{(\E[N]-256)^2}
\le
\frac{\E[N]}{(\E[N]-256)^2}
\le \frac{1024}{768^2}
<\frac1{16}.
\]

\smallskip
\noindent\emph{Step 4: a Sperner bound for the symmetrized sum.}
Let
\[
U:=\sum_{j=1}^M Y_j.
\]
We claim that for every $s\in\R$,
\begin{equation}\label{eq:symm-smallball}
\Prob(\abs{U-s}\le 2r_0)\le \frac18.
\end{equation}
Indeed,
\[
\Prob(\abs{U-s}\le 2r_0)
\le
\Prob(\abs{U-s}\le 2r_0,\ N\ge256)+\Prob(N<256).
\]
We have already bounded the second term by $1/16$, so it remains to estimate
the first one.

Condition on the absolute values $(\abs{Y_j})_{j=1}^M$.  Since each $Y_j$ is
symmetric, conditional on $(\abs{Y_j})_j$ the signs of the nonzero $Y_j$ are
independent Rademacher variables.  On the event $\{N\ge256\}$ let
\[
J:=\{j:\abs{Y_j}\ge a_0\},
\qquad
m:=\#J\ge256.
\]
After freezing the inactive coordinates and the magnitudes $\abs{Y_j}$, the
condition $\abs{U-s}\le 2r_0$ becomes
\[
\sum_{j\in J}\varepsilon_j w_j \in I,
\]
where $w_j:=\abs{Y_j}\ge a_0$ and $I$ is an interval of diameter
$4r_0=a_0$.  Since $a_0<2w_j$ for every $j\in J$, the family of sign patterns
for which the weighted sum falls in $I$ is an antichain in $\{\pm1\}^J$.
Indeed, if two such sign patterns were comparable in the natural subset order,
their sums would differ by at least $2w_j\ge2a_0>a_0$, which is impossible
inside an interval of diameter $a_0$.  By Sperner's theorem, the conditional
probability of this event is therefore at most
\[
\frac{\binom{m}{\lfloor m/2\rfloor}}{2^m}
\le \frac1{\sqrt m}
\le \frac1{16},
\]
because $m\ge256$.  Averaging over the conditioning yields
\[
\Prob(\abs{U-s}\le 2r_0,\ N\ge256)\le \frac1{16}.
\]
Together with the bound for $\Prob(N<256)$, this proves
\eqref{eq:symm-smallball}.

\smallskip
\noindent\emph{Step 5: unsymmetrization.}
Let
\[
R:=\sum_{i\notin G_1\cup\cdots\cup G_M} e_iD_i.
\]
Then $R$ is independent of $U$, since it is built from coordinates disjoint
from those entering the block sums.  Also,
\[
X_e(\omega_1)-X_e(\omega_2)=U+R.
\]
Fix $t\in\R$.  If both $\abs{X_e(\omega_1)-t}\le r_0$ and
$\abs{X_e(\omega_2)-t}\le r_0$, then
\[
\abs{X_e(\omega_1)-X_e(\omega_2)}\le 2r_0.
\]
Hence
\[
\Prob(\abs{X_e-t}\le r_0)^2
\le
\Prob(\abs{X_e(\omega_1)-X_e(\omega_2)}\le 2r_0)
=
\Prob(\abs{U+R}\le 2r_0).
\]
Conditioning on $R$ and using \eqref{eq:symm-smallball},
\[
\Prob(\abs{U+R}\le 2r_0)
\le
\sup_{s\in\R}\Prob(\abs{U-s}\le 2r_0)
\le \frac18.
\]
Therefore
\[
\Prob(\abs{X_e-t}\le r_0)\le \frac1{\sqrt8}<\frac38.
\]
This proves \eqref{eq:onevar-smallball}, and with it the lemma.
\end{proof}

Once this shifted anticoncentration estimate is available, the lower bound for
$\eps$ follows from the same factorization of $X^2-Y^2$ that was
already used in the compactness proof.

\begin{proposition}[Diffuse branch]\label{prop:diffuse}
If the diffuse alternative of Lemma~\ref{lem:dichotomy} holds, then
\[
\eps\ge c_{\mathrm{cross}}(\delta)>0.
\]
\end{proposition}

\begin{proof}
Let
\[
h_\pm:=X_{(a\pm b)\mathbf 1_H}.
\]
Then
\[
X^2-Y^2
=2\bigl(X_d+h_+/\sqrt2\bigr)\bigl(X_c+h_-/\sqrt2\bigr),
\]
and therefore
\[
\abs{X^2-Y^2}
\ge
2\min\Bigl(\abs{X_d+h_+/\sqrt2},\abs{X_c+h_-/\sqrt2}\Bigr)^2.
\]
Condition on the head variables $h_\pm$ and apply Lemma~\ref{lem:cross} with
the shifts $u=h_-/\sqrt2$ and $v=h_+/\sqrt2$.  This yields a uniform positive
lower bound for the conditional expectation of the squared minimum, hence for
$\eps$.  The resulting constant depends only on $\delta$.
\end{proof}

\subsection{Conclusion of the second proof}

At this point the proof is immediate: the dichotomy shows that every
orthogonal pair falls into exactly one of the two branches, and each branch
gives a quantitative lower bound for the same quantity $\eps$.

\begin{proof}[Second proof of Theorem~\ref{thm:main}]
By Proposition~\ref{prop:orth-reduction}, it is enough to prove a uniform
orthogonal lower bound.  Let $a,b$ be orthogonal unit vectors and form $X$
and $Y$ as above.  Apply the dichotomy of Lemma~\ref{lem:dichotomy}.

If the concentrated alternative holds, then
Proposition~\ref{prop:concentrated} yields
\[
\eps\ge c_{\mathrm{line}}(\delta,\eta).
\]
If the diffuse alternative holds, then Proposition~\ref{prop:diffuse} yields
\[
\eps\ge c_{\mathrm{cross}}(\delta).
\]
Since $\eps\le 2\Delta$, we conclude that
\[
\Delta
=\norm{|X|-|Y|}_{L^2}
\ge
\frac12\min\bigl(c_{\mathrm{line}}(\delta,\eta),c_{\mathrm{cross}}(\delta)\bigr)
=:c_{\delta,\eta}>0.
\]
The orthogonal reduction then gives
\[
\min_{\sigma\in\{\pm1\}}\norm{X-\sigma Y}_{L^2}
\le \frac{\sqrt2}{c_{\delta,\eta}}\norm{|X|-|Y|}_{L^2}
\]
for all $X,Y\in \overline{\Span}(\eta,\xi_i:i\ge2)$, completing the proof of
Theorem~\ref{thm:main}.
\end{proof}

\begin{corollary}[Finite-dimensional one-exception stable phase retrieval]
\label{cor:finite-one-exception}
Let $n\ge2$.  Then there exists $C_{\delta,\eta}\ge1$ such that, for all
\[
X,Y\in \Span(\eta,\xi_2,\dots,\xi_n),
\]
one has
\[
\min_{\sigma\in\{\pm1\}}\norm{X-\sigma Y}_{L^2}
\le
C_{\delta,\eta}\norm{|X|-|Y|}_{L^2}.
\]
\end{corollary}

\begin{proof}
This is an immediate consequence of Theorem~\ref{thm:main}, since
\[
\Span(\eta,\xi_2,\dots,\xi_n)\subset \overline{\Span}(\eta,\xi_i:i\ge2).
\]
\end{proof}

\section{Comments and Open Problems}\label{sec:comments}

The compactness and quantitative proofs suggest several natural directions for
further work.  Although the theorem proved here settles the qualitative
one-exception problem in the real Hilbertian setting, the mechanisms that enter
the proof are flexible enough that one is led almost immediately to adjacent
questions.  Some of these concern extensions of the model, while others concern
the sharp quantitative form of the present result.  We briefly record several
of them here.

\begin{remark}[Complex-valued models]
The complex one-exception theorem already exhibits the same head--tail
architecture, but the endpoint geometry is richer because the analogue of the
phase-retrieval cone is foliated by complex lines and conjugate planes rather
than by two real lines.  The role of the upper $L^1$ bound is then supplemented
by a generalized-circle gap.  The real argument presented here should be viewed
as the cleanest model case for that geometry.  This makes the complex setting
an especially interesting next test case: it is close enough to the present
problem that one expects the same structural ideas to survive, but different
enough that the endpoint classification becomes genuinely more delicate.
\end{remark}

\begin{remark}[More general symmetries]
Classical phase retrieval is recovery modulo the group $\{\pm1\}$.  It is
natural to ask for analogous theorems modulo larger compact symmetry groups or
for quotient metrics arising from finite reflection groups.  The compactness
proof is particularly well suited to that direction because it separates the
problem into an exact limit identity and a geometric classification of the
corresponding support set.  For that reason, this is not merely a formal
generalization: it points toward a broader geometric framework in which the
support theorem, rather than the specific real-sign symmetry, is the central
object.
\end{remark}

The next questions concern the extent to which the present theorem can be
sharpened.  The first asks for a genuinely effective version of the one-exception
theorem itself.

\begin{question}
Can one obtain a fully quantitative version of Theorem~\ref{thm:main} with one
exceptional coordinate, with constants depending explicitly on
$\delta$ and quantitative parameters of the exceptional variable $\eta$?
\end{question}

This question is particularly interesting because the two proofs in the paper
point in different directions: the compactness proof is robust but soft, while
the second proof isolates several places where explicit constants might, at
least in principle, be tracked.

Closely related to that issue is the problem of identifying the correct endpoint
assumption on the good coordinates.

\begin{question}
What is the optimal endpoint replacement for the upper bound
$\norm{\xi_i}_{L^1}\le1-\delta$?  In the real case this bound excludes the
Rademacher obstruction, but it is not obvious whether a weaker quantitative
non-flatness parameter could suffice.
\end{question}

This is again an interesting structural question rather than a matter of
constant optimization.  The present upper bound is easy to interpret and
natural for the proof, but it would be valuable to know whether the real
obstruction is exactly Rademacher degeneracy or some larger non-flatness
phenomenon.

Finally, one may ask how much of the theorem belongs specifically to Hilbert
space geometry, and how much should persist in a wider Banach-function-space
setting.

\begin{question}
Which Banach-function-space analogues of the present theorem remain true once
the ambient space is no longer Hilbertian?  The orthogonal reduction is
special to Hilbert spaces, so a different mechanism would be needed.
\end{question}

This problem seems especially compelling because the function-space formulation
of stable phase retrieval is not inherently Hilbertian, even though the present
proof is.  Any satisfactory extension would likely require replacing
orthogonality by a more flexible geometric notion while retaining enough
structure to control the modulus map.

\begin{remark}
The quantitative proof is stronger whenever one seeks effective constants, but
the compactness proof is likely the more durable template for extensions.  It
uses only bounded probes, a profile decomposition, and a structural
description of diffuse limits.  Those ingredients are robust enough to adapt
to complex models, to settings with exceptional coordinates, and plausibly to
retrieval modulo more complicated symmetries.  Taken together, these questions
suggest that the present theorem is less an endpoint than a model case: it
captures a setting in which the geometry can still be classified cleanly, while
already displaying most of the phenomena that one would expect to confront in
more general phase-retrieval problems.
\end{remark}

\nocite{calderbank2022stable}
\nocite{bandeira2014saving}
\nocite{eldar2014stability}
\nocite{cahill2016infinite}
\nocite{alaifari2017continuous}
\nocite{alaifari2021gabor}
\nocite{grohs2020survey}
\nocite{FOTP}
\nocite{christ2022examples}
\nocite{localphase}
\nocite{chen2019phase}
\nocite{MR3746047}
\nocite{KS}
\nocite{DB}
\nocite{MR3260258}
\nocite{MR3069958}
\nocite{MR4298599}
\nocite{sanz1984phase}
\nocite{balan2013}
\nocite{bartusel2021phase}
\nocite{casazza2017weak}
\nocite{casazza2017norm}
\nocite{casazza2013}
\nocite{AG-continuous-frames}
\nocite{fg22}
\nocite{MR3901674}
\nocite{grohs2021stable}
\nocite{grohs20212}
\nocite{luef2021gaussian}
\nocite{kashin2021sampling}
\nocite{LT21}
\nocite{DPSTT}
\nocite{NOU}
\nocite{akrami2021note}
\nocite{MR3367650}
\nocite{MR3275625}
\nocite{CCS-JMAA}
\nocite{CDOS}
\nocite{CHL}
\nocite{FG-atomic}
\nocite{MR2011364}
\nocite{MR0423039}
\nocite{our-paper}

\bibliographystyle{plain}
\bibliography{refs}

\end{document}
